\documentclass[a4paper,11pt]{article}
\usepackage[utf8]{inputenc}
\usepackage[T1]{fontenc}
\usepackage[french]{babel}
\usepackage[left=1.8cm,right=1.8cm,top=1.5cm,bottom=1.5cm]{geometry}
\usepackage{xcolor}
\usepackage{hyperref}
\usepackage{fontawesome5}
\usepackage{parskip}

% --- Couleurs (Vert Valeo) ---
\definecolor{primary}{RGB}{0, 128, 0}
\hypersetup{colorlinks=true, urlcolor=primary}

\begin{document}

% --- EXPÉDITEUR ---
\noindent
\begin{minipage}[t]{0.5\textwidth}
    \textbf{Loïc Aron MBASSI EWOLO}\\
    Élève-Ingénieur Bac+5 -- Double diplôme ENSEA\\[3pt]
    \faMapMarker\ Cergy, France\\
    \faPhone\ (+33) 7 59 52 33 98\\
    \faEnvelope\ \href{mailto:loicaron.mbassiewolo@ensea.fr}{loicaron.mbassiewolo@ensea.fr}
\end{minipage}
\hfill
% --- DATE ---
\begin{minipage}[t]{0.4\textwidth}
    \raggedleft
    Cergy, le 27 février 2026
\end{minipage}

\vspace{1.2cm}

% --- DESTINATAIRE ---
\hfill
\begin{minipage}[t]{0.5\textwidth}
    \raggedleft
    \textbf{Valeo -- Division Light}\\
    Service Recrutement\\
    Bobigny, France
\end{minipage}

\vspace{1cm}

\noindent\textbf{Objet :} Candidature -- Stagiaire Ingénieur Data Science \& Intelligence Artificielle, SDV Lighting

\vspace{0.8cm}

Madame, Monsieur,

Entraîner un modèle de deep learning jusqu'à convergence, puis l'optimiser pour du hardware embarqué contraint : c'est le fil conducteur de mes projets, du laboratoire MILA (Institut Québécois d'IA) au challenge CoHoMa de robotique autonome à l'ENSEA. En double diplôme ENSEA/Polytechnique Yaoundé, spécialisé en IA et systèmes embarqués, je souhaite contribuer au développement de la fonction anti-éblouissement par IA au sein de l'équipe SDV Lighting de Valeo.

Au MILA, j'ai étudié le phénomène de Grokking sous PyTorch en reproduisant des expériences état de l'art, avec une analyse mathématique de la convergence des modèles, compétences directement applicables à l'entraînement et l'évaluation des algorithmes anti-glare. Mon projet UROP (mentors Google DeepMind) portait sur la classification d'anomalies cardiaques à partir d'images d'ECG via TensorFlow/Keras, combinant traitement du signal et Computer Vision, un pipeline transférable à l'analyse d'images pour la détection d'éblouissement.

Sur le volet embarqué, j'optimise actuellement des modèles YOLOv10 sur Nvidia Jetson (challenge CoHoMa, ENSEA) sous contraintes CPU/mémoire, et mon stage au laboratoire IoT (ENSPY) m'a formé au déploiement TinyML sur Raspberry Pi avec TensorFlow Lite. Ces expériences répondent à l'enjeu de compatibilité embarquée de vos algorithmes. Côté LLM, j'ai conçu un système multi-agents utilisant RAG et Gemini via Google ADK. Mon projet de Fault Detection sur drone sous MATLAB/Simulink et le hackathon IndabaX 2025 (1\textsuperscript{ère} place), où j'ai construit des pipelines de données et tableaux de bord Streamlit, complètent ce profil.

Je suis disponible dès avril 2026 pour un stage de 4 à 6 mois. Contribuer à la fonction anti-glare chez Valeo Lighting me permettrait de combiner data science, optimisation de modèles et déploiement embarqué sur un projet à impact concret pour la sécurité routière. Mon anglais professionnel, acquis lors de collaborations internationales (MILA au Canada, SOSDOCTA en Belgique), me permettra d'évoluer dans votre environnement multiculturel.

Je reste à votre disposition pour un entretien afin de discuter de ma candidature.

Je vous prie d'agréer, Madame, Monsieur, l'expression de mes salutations distinguées.

\vspace{0.8cm}

\textbf{Loïc Aron MBASSI EWOLO}\\
\textit{Élève-Ingénieur ENSEA / Polytechnique Yaoundé}

\end{document}
